\chapter{Related Work}
\label{ch:related_work}

This chapter reviews prior research relevant to the challenges this thesis addresses. First, we examine work on estimating musical difficulty from symbolic score representations, which establishes that symbolic notation carries usable difficulty signal but also exposes recurring weaknesses in label provenance and model transparency. Second, we review the argument that difficulty is a multi-dimensional quantity, which directly motivates the organization of our descriptor set. Third, we describe the RubricNet architecture, whose additive structure forms the methodological core of this thesis. Fourth, we review fuzzy rule-based approaches to interpretable music classification, from which two additional baselines used in the evaluation are drawn. Finally, we survey guitar-specific playability metrics and the datasets from which our corpus is assembled.

\section{Symbolic Difficulty Estimation}
Automatic difficulty assessment from musical scores is a small but growing subfield of music information retrieval. Yaroch~\cite{yaroch2025measure} estimated piano score difficulty on a measure-by-measure basis from MusicXML, extracting symbolic features such as note density, interval sizes, and basic rhythmic properties, and training linear regression models and convolutional neural networks on them. The study demonstrated that symbolic notation alone carries usable difficulty signal, but it also illustrates two recurring weaknesses of the area: the difficulty labels were self-assigned rather than drawn from an established grading system, and the better-performing models were black boxes that could not explain their predictions. Both issues are consequential in a pedagogical setting, where a difficulty estimate is only actionable if a teacher can trust its provenance and verify its reasoning. This thesis addresses the first weakness by grading against an established external repertoire database, and the second by adopting an interpretable-by-construction architecture.

\section{Multi-Dimensional Difficulty Representation}
Difficulty is not a single quantity. A piece can be hard to read but easy to execute, or technically demanding in one hand while trivial in the other. Ramoneda et al.~\cite{ramoneda2024combining} make this argument explicit for piano repertoire, decomposing difficulty into notation (reading complexity), technique (physical execution), and expressiveness (interpretative demands), and showing that estimators trained on the separate dimensions and then combined outperform a flat, single-vector feature representation. This finding directly shapes the descriptor design in this thesis: rather than pooling all features, the guitar descriptors are organized into left-hand technique, right-hand technique, and global score complexity, and Chapter~\ref{ch:evaluation} tests whether the combination in fact outperforms each dimension in isolation.

\section{The RubricNet Model}
The methodological core of this thesis is RubricNet, introduced by Ramoneda et al.~\cite{ramoneda2024rubricnet} as a white-box ordinal classifier for piano difficulty on the CIPI dataset. Its defining property is an additive structure: each hand-crafted descriptor passes through its own small subnetwork, and the scalar outputs are summed into a single difficulty score,
\begin{equation}
    S(x) = \sum_{i=1}^{M} g_i(x_i),
\end{equation}
where $x_i$ is the $i$-th descriptor and $g_i$ its subnetwork. Because no descriptor interacts with any other before the summation, the contribution of each one to the final score can be read off directly from the trained model. The design deliberately mirrors a pedagogical rubric, in which a teacher scores independent aspects of a performance and adds the scores together.

RubricNet achieved accuracy competitive with black-box baselines on piano data while remaining fully transparent. Its descriptors, however, are keyboard-specific --- separate left- and right-hand staves, key spans, piano fingering patterns --- and cannot be applied to guitar scores, where a single staff encodes the work of both hands and the same pitch can be played at several fretboard positions. Adapting the architecture therefore requires redesigning the entire descriptor layer, which constitutes the main technical work of this thesis. To the best of our knowledge, this is the first application of the additive rubric architecture outside piano repertoire.

\section{Fuzzy Rule-Based Music Classification}
\label{sec:rw_fuzzy}
The additive rubric is not the only route to interpretable music classification. A separate line of work builds classifiers from \emph{fuzzy rules}~\cite{zadeh1965fuzzy} --- natural-language statements of the form ``if fret entropy is very high, then the piece is hard'' --- whose interpretability lies in the rules themselves rather than in a decomposition of a numeric score. Vatolkin and Rudolph~\cite{vatolkin2015interpretable} introduced this approach to music categorization, inducing primitive rules over high-level audio features by an exhaustive relevance-ranked search and showing that the resulting rule bases remain readable by non-experts. Heerde et al.~\cite{heerde2020comparing} subsequently compared three fuzzy rule-based approaches on music genre classification: the complete search over primitive rules, an evolutionary optimization of entire rule bases, and fuzzy pattern trees~\cite{huang2008patterntrees, senge2011topdown, senge2015fast}, which compose fuzzy statements with logical operators into per-class trees. All three achieved usable genre accuracy while producing rules a listener can inspect directly.

This line of work matters to this thesis for two reasons. First, it offers a second, independently developed paradigm of interpretability against which RubricNet can be compared: readable rules over named features, rather than summed per-descriptor scores. Second, the methods are class-agnostic --- they treat labels as unordered categories --- which makes them a natural instrument for isolating the value of RubricNet's ordinal design. Section~\ref{sec:fuzzy_method} adapts two of the three approaches (the complete search and fuzzy pattern trees; the evolutionary variant is excluded as nondeterministic and, in the source comparison, the least readable) from genre classification to difficulty grading, where to our knowledge they have not previously been applied.

\section{Guitar Playability Metrics}
For the guitar specifically, V{\'a}squez et al.~\cite{vasquez2023quantifying} proposed a playability metric for rhythm-guitar chord progressions, built from expert-informed criteria: uncommonness of chord shapes, finger displacement between chords, barre usage, right-hand strumming complexity, and rhythmic regularity. Their work is an important precedent for descriptor design --- several of the descriptors developed in this thesis, such as the barre ratio, chord stretch, and shift distance, have direct analogues in their rubric --- but its scope is limited to strummed chord charts in popular music. Solo classical repertoire is a substantially harder domain: textures are polyphonic and contrapuntal, the left hand moves horizontally through fretboard positions rather than between a fixed vocabulary of chord shapes, and the right hand executes arpeggiation and voice-separation patterns that have no counterpart in strumming. This thesis extends the playability-descriptor idea to that domain and connects it, for the first time, to an external professionally graded repertoire resource.

\section{Datasets}
The dataset constructed in Chapter~\ref{ch:method} draws on several resources that play distinct roles. Difficulty labels come from GuitarBurst~\cite{guitarburst}, an online graded repertoire database whose 1--20 difficulty grades serve as the prediction target. Symbolic scores come from two published research datasets: DadaGP~\cite{sarmento2021dadagp}, a large corpus of GuitarPro tablature transcriptions in a tokenized format, from which the classical guitar subset is used, and GAPS~\cite{riley2024gaps}, a smaller, high-quality classical guitar dataset with MusicXML scores aligned to audio performances. The bulk of the symbolic scores, however, come from ClassClef, an online classical-guitar score library rather than a research dataset; because it has no associated publication, it is described where the corpus is actually assembled (Chapter~\ref{ch:method}) rather than surveyed here. None of these resources alone pairs symbolic scores with difficulty grades at useful scale, which is why the matching and alignment pipeline described in Chapter~\ref{ch:method} was necessary.
